% !TeX root = ../../paper.tex

We built a VR horror experience in which every dreadful thing the player must touch is really there to be touched, in proxy form: rubber balls for eyes, gummy worms for worms, a hidden water spray for a burst pipe, synthetic webbing for cobwebs. The technical recipe is deliberately reproducible---consumer hardware, asset-store content, spatial anchors, printed fiducial markers, and props from a supermarket and a Halloween aisle. Our within-subject comparison of the experience with and without these proxies, measured with the IPQ, \todoin{summarize the outcome in one sentence}.

Beyond the specific numbers, the project's practical lesson is that the hard part of passive haptics is not the props but the alignment: the illusion lives or dies on whether the hand meets the object exactly where the eyes expect it, which is a calibration problem more than a design problem. Future work has an obvious next step in restoring the measures we cut---SAM and physiological data such as heart rate variability would separate presence from fright, which the IPQ alone cannot do. A per-challenge manipulation could then identify which kinds of touch (immersion, sudden stimulus, clinging texture) carry the effect, and whether dread needs the hand to be met at all, or only the fear that it might be.
